\begin{abstract}
本稿では，画像中の物体らしい領域を検出するタスクと，その領域に対応するカテゴリをアノテーションするタスクを明確に分離した物体アノテーション手法を提案する．
物体領域の検出にはカテゴリに依存しないDEtection TRansformer（DETR）を用いて1回の推論でバウンディングボックスを得る．
カテゴリのアノテーションには，各領域のクロップに対する視覚言語モデル（Vision-Language Model，VLM）のゼロショット分類を用いる．
両者が分離されているため，一度検出した領域に対して候補語彙集合を変更するだけでカテゴリのアノテーションを繰り返すことができる．
さらに，WordNetの上位語・下位語関係をたどることで，追加学習なしに「animal → dog → doberman」のような粒度可変のアノテーションを実現する．
COCO 2017 Validationの500枚を用いた評価により，以下を確認した．
第一に，ImageNetで学習した閉域の分類モデルに比べてContrastive Language-Image Pre-training（CLIP）は画像単位の単語推定精度が統計的に有意に高く，カテゴリのアノテーションにはオープンボキャブラリなVLMが必要である．
第二に，クラス非依存の位置検出精度において，DETRはOWL-ViTおよびGrounding DINOに劣らない．
第三に，提案パイプライン全体のカテゴリアノテーション精度はTop-1正解率0.604，Top-5正解率0.894であり，クロップを正方形にパディングしてVLMへ入力することでTop-1正解率が約0.08向上する．
第四に，WordNet階層の最初の細分化の段階で確信度が0.642から0.442へ大きく低下する．
また，下位語が1語しかない分岐が細分化の段階全体の約18\%で生じ，これが深い階層における見かけの確信度を押し上げていることを明らかにした．
さらに，COCOの80クラスで学習したYOLOと同一の語彙で比較し，80クラスの範囲内ではYOLOのTop-1正解率が0.98と提案手法を大きく上回るものの，その差は検出器ではなく分類器に起因すること，およびYOLOはCOCO外の物体をbirdやdog等の80クラス内の語でしかアノテーションできないことを示した．
加えて，深さ0のカテゴリをYOLOの予測とし，深さ1以降の細分化のみをVLMとWordNetで行うハイブリッド構成により，深さ0のTop-1正解率0.981を保ったまま，粒度可変のアノテーションの対象となるインスタンスを提案手法の1.65倍に増やせること，およびVLMの確信度が0.5以上の場合のみCOCO外の語に切り替えることで，\add{COCOの精度をほぼ損なわずに}，COCO外の物体の実画像にも対応できることを示した．
\add{さらに，別の500枚とLVISの細かい粒度の正解を用いて，これらの結論が別の画像でも成り立つこと，深さ0が正しくても深さ1の推定が正しいのは半数未満であること，および累積確信度による細分化の停止によって階層的適合率が0.594から0.801へ向上することを示した．}
COCOに含まれない語を主要な被写体とする実画像に対しても追加学習なしに正しくアノテーションできる一方，空を背景とした塔状の構造物ではDETRの検出自体が失敗することも確認した．
\end{abstract}

\noindent
{\bf キーワード：}オープンボキャブラリ物体アノテーション，DETR，視覚言語モデル，CLIP，WordNet，粒度可変アノテーション
